\documentclass[12pt,a4paper]{article}
\usepackage{fontspec}
\setmainfont{TeX Gyre Termes}
\usepackage[top=1.6cm,bottom=1.6cm,left=1.8cm,right=1.5cm]{geometry}
\usepackage{amsmath,amssymb}
\usepackage{enumitem}
\usepackage{multicol}
\usepackage{xcolor}
\usepackage{array}
\usepackage{titlesec}
\setlength{\parindent}{0pt}
\setlist{nosep}
\definecolor{ddo}{HTML}{C0392B}
\definecolor{dxanh}{HTML}{1F4E79}

% ===== Cờ hiển thị đáp án: mặc định HIỆN (bản GV) =====
\newif\ifANS \ANStrue
\ifdefined\HIDEANS \ANSfalse \fi

\newcommand{\dapan}[1]{\ifANS\textbf{\color{ddo}#1}\else#1\fi}
\newcommand{\loigiai}[1]{\ifANS\par\smallskip{\color{dxanh}\textbf{Lời giải. }#1}\par\fi}
\newcommand{\dapanngan}[1]{\ifANS~~\textbf{\color{ddo}[Đáp án: #1]}\fi}
% ý đúng/sai: #1 = Đ hoặc S, #2 = nội dung
\newcommand{\yt}[2]{\item #2\ifANS~~\textbf{\color{ddo}(#1)}\fi}

\titleformat{\section}{\large\bfseries\color{dxanh}}{}{0pt}{}
\titlespacing{\section}{0pt}{8pt}{4pt}

\newcommand{\cau}{\medskip\textbf{Câu \refstepcounter{cau}\thecau.}~}
\newcounter{cau}
\newcommand{\resetcau}{\setcounter{cau}{0}}

\begin{document}

% ====== TIÊU ĐỀ ======
\begin{center}
{\large\bfseries ĐỀ ÔN TẬP HÈ -- KHOA HỌC TỰ NHIÊN 6}\\[2pt]
{\itshape (Dành cho học sinh lớp 5 lên lớp 6) -- Năm học 2025--2026}\\[4pt]
{\bfseries Thời gian làm bài: 45 phút \quad | \quad Mã đề: 101}
\ifANS\\[2pt]{\color{ddo}\bfseries (BẢN GIÁO VIÊN -- CÓ ĐÁP ÁN \& LỜI GIẢI)}\fi
\end{center}
\hrule
\smallskip
Họ và tên: \dotfill\quad Lớp: \rule{2.2cm}{0.4pt}
\smallskip
\hrule

% ====== PHẦN I ======
\section{PHẦN I. TRẮC NGHIỆM NHIỀU LỰA CHỌN (mỗi câu chọn một phương án)}
\resetcau

\cau Khí nào sau đây chiếm khoảng $78\%$ thể tích không khí?
\begin{multicols}{2}
\begin{itemize}[label={},leftmargin=1.2em]
\item A. Oxygen
\item B. Carbon dioxide
\item \dapan{C. Nitrogen}
\item D. Hơi nước
\end{itemize}
\end{multicols}
\loigiai{Nitrogen chiếm khoảng $78\%$ thể tích không khí, oxygen khoảng $21\%$, phần còn lại là carbon dioxide, hơi nước và một số khí khác.}

\cau Phơi quần áo ướt ngoài nắng, sau vài giờ quần áo khô. Nước trong quần áo đã biến đổi như thế nào?
\begin{multicols}{2}
\begin{itemize}[label={},leftmargin=1.2em]
\item A. Nước đông đặc
\item B. Nước ngưng tụ
\item \dapan{C. Nước bay hơi vào không khí}
\item D. Nước thấm ngược vào vải
\end{itemize}
\end{multicols}
\loigiai{Dưới ánh nắng, nước (thể lỏng) nhận nhiệt chuyển thành hơi nước (thể khí) bay vào không khí -- đó là sự bay hơi, nên quần áo khô.}

\cau Ở vùng núi cao, nước thường sôi ở nhiệt độ thấp hơn $100^\circ$C. Nguyên nhân chủ yếu là
\begin{itemize}[label={},leftmargin=1.2em]
\item A. nước ở vùng núi không tinh khiết
\item B. không khí trên núi lạnh hơn
\item \dapan{C. áp suất khí quyển trên cao thấp hơn nên nhiệt độ sôi của nước giảm}
\item D. nước trên núi chảy nhanh hơn
\end{itemize}
\loigiai{Càng lên cao áp suất khí quyển càng giảm, làm nước sôi ở nhiệt độ thấp hơn $100^\circ$C.}

\cau Trường hợp nào sau đây là một dung dịch?
\begin{multicols}{2}
\begin{itemize}[label={},leftmargin=1.2em]
\item A. Cát khuấy đều trong nước
\item B. Dầu ăn khuấy với nước
\item \dapan{C. Muối ăn tan hoàn toàn trong nước}
\item D. Gạo trộn lẫn với đỗ xanh
\end{itemize}
\end{multicols}
\loigiai{Muối tan hoàn toàn và phân bố đều trong nước tạo thành dung dịch (nước muối). Cát, dầu ăn không tan; gạo và đỗ chỉ là hỗn hợp.}

\cau Để tách muối ăn ra khỏi dung dịch nước muối, ta nên
\begin{itemize}[label={},leftmargin=1.2em]
\item A. dùng nam châm hút muối
\item B. dùng giấy lọc để lọc muối
\item \dapan{C. đun nóng dung dịch cho nước bay hơi, muối kết tinh ở lại}
\item D. để yên cho muối tự lắng xuống
\end{itemize}
\loigiai{Muối tan trong nước nên không lọc hay lắng được. Đun nóng làm nước bay hơi, muối không bay hơi sẽ kết tinh và ở lại.}

\cau Hỗn hợp nào sau đây có thể tách bằng nam châm?
\begin{multicols}{2}
\begin{itemize}[label={},leftmargin=1.2em]
\item A. Muối và nước
\item B. Đường và nước
\item \dapan{C. Mạt sắt và cát}
\item D. Dầu ăn và nước
\end{itemize}
\end{multicols}
\loigiai{Sắt bị nam châm hút còn cát thì không, nên dùng nam châm tách được mạt sắt khỏi cát.}

\cau Chất ở trạng thái khí có đặc điểm nào sau đây?
\begin{itemize}[label={},leftmargin=1.2em]
\item A. Có hình dạng và thể tích xác định
\item B. Có hình dạng xác định nhưng thể tích thay đổi
\item \dapan{C. Không có hình dạng xác định, chiếm toàn bộ không gian của vật chứa}
\item D. Chỉ chiếm phần đáy của vật chứa
\end{itemize}
\loigiai{Chất khí không có hình dạng và thể tích xác định; nó chiếm đầy toàn bộ không gian vật chứa.}

\cau Khi đun nóng, một mẩu nến rắn chảy thành nến lỏng. Đây là sự chuyển trạng thái
\begin{multicols}{2}
\begin{itemize}[label={},leftmargin=1.2em]
\item A. lỏng sang rắn
\item \dapan{B. rắn sang lỏng}
\item C. rắn sang khí
\item D. lỏng sang khí
\end{itemize}
\end{multicols}
\loigiai{Nến từ thể rắn nóng chảy sang thể lỏng -- chuyển trạng thái rắn sang lỏng, không tạo chất mới.}

\cau Sự biến đổi hóa học là sự biến đổi
\begin{itemize}[label={},leftmargin=1.2em]
\item A. chỉ làm thay đổi hình dạng của chất
\item B. chỉ làm thay đổi trạng thái của chất
\item \dapan{C. có sự tạo thành chất mới}
\item D. làm cho chất bị chia nhỏ ra
\end{itemize}
\loigiai{Dấu hiệu cơ bản của biến đổi hóa học là tạo thành chất mới (thường đổi màu, mùi, tính tan).}

\cau Hiện tượng nào sau đây là một biến đổi hóa học?
\begin{itemize}[label={},leftmargin=1.2em]
\item A. Nước đá tan thành nước
\item B. Bẻ đôi một viên phấn
\item \dapan{C. Đinh sắt để lâu trong không khí ẩm bị gỉ}
\item D. Hòa tan đường vào nước
\end{itemize}
\loigiai{Đinh sắt bị gỉ tạo ra chất mới (gỉ sắt) màu nâu đỏ, giòn -- là biến đổi hóa học. Ba hiện tượng còn lại chỉ đổi trạng thái/hình dạng.}

\cau Dụng cụ nào sau đây dùng để đo chính xác thể tích chất lỏng?
\begin{multicols}{2}
\begin{itemize}[label={},leftmargin=1.2em]
\item A. Cốc thủy tinh
\item \dapan{B. Ống đong}
\item C. Đũa thủy tinh
\item D. Kẹp ống nghiệm
\end{itemize}
\end{multicols}
\loigiai{Ống đong có vạch chia thể tích nên dùng để đo chính xác thể tích chất lỏng.}

\cau Việc làm nào sau đây KHÔNG đúng quy tắc an toàn trong phòng thực hành?
\begin{itemize}[label={},leftmargin=1.2em]
\item A. Đeo kính bảo hộ khi làm thí nghiệm
\item \dapan{B. Nếm thử hóa chất để biết đó là chất gì}
\item C. Đọc kỹ nhãn hóa chất trước khi dùng
\item D. Rửa tay sạch bằng xà phòng sau khi thực hành
\end{itemize}
\loigiai{Tuyệt đối không nếm hoặc ngửi trực tiếp hóa chất vì có thể gây ngộ độc, bỏng.}

% ====== PHẦN II ======
\section{PHẦN II. TRẮC NGHIỆM ĐÚNG/SAI (mỗi ý chọn Đúng hoặc Sai)}
\resetcau

\cau Cho các phát biểu sau về nước, không khí và đất:
\begin{enumerate}[label=\alph*),leftmargin=1.6em]
\yt{Đ}{Nước có thể tồn tại ở cả ba thể: rắn, lỏng và khí.}
\yt{S}{Không khí có màu xanh nhạt và có vị ngọt.}
\yt{Đ}{Mùn trong đất cung cấp chất dinh dưỡng cho cây trồng.}
\yt{S}{Nitrogen chiếm khoảng $21\%$ thể tích không khí.}
\end{enumerate}
\loigiai{a) Đúng -- nước có ba thể. b) Sai -- không khí trong suốt, không màu, không mùi, không vị. c) Đúng -- mùn giàu chất hữu cơ, dinh dưỡng. d) Sai -- nitrogen chiếm $\approx 78\%$, oxygen mới $\approx 21\%$.}

\cau Cho các phát biểu sau về sự biến đổi của chất:
\begin{enumerate}[label=\alph*),leftmargin=1.6em]
\yt{Đ}{Đốt cháy một tờ giấy thành tro là biến đổi hóa học.}
\yt{S}{Nước đá tan thành nước lỏng là biến đổi hóa học.}
\yt{Đ}{Có thể nhận biết biến đổi hóa học qua dấu hiệu đổi màu, đổi mùi hoặc đổi tính tan.}
\yt{Đ}{Thức ăn để lâu bị ôi thiu là một biến đổi hóa học.}
\end{enumerate}
\loigiai{a) Đúng -- tạo ra tro là chất mới. b) Sai -- chỉ đổi trạng thái (biến đổi vật lí). c) Đúng. d) Đúng -- ôi thiu sinh chất mới có mùi, vị khác.}

% ====== PHẦN III ======
\section{PHẦN III. TRẢ LỜI NGẮN}
\resetcau

\cau Ở điều kiện thường, nước sôi ở nhiệt độ bao nhiêu độ C? \dapanngan{$100$}
\loigiai{Ở điều kiện thường, nước sôi ở $100^\circ$C.}

\cau Nitrogen chiếm khoảng bao nhiêu phần trăm thể tích không khí? \dapanngan{$78$}
\loigiai{Nitrogen chiếm khoảng $78\%$ thể tích không khí.}

\cau Ở điều kiện thường, nước đông đặc (hóa thành nước đá) ở nhiệt độ bao nhiêu độ C? \dapanngan{$0$}
\loigiai{Ở điều kiện thường, nước đông đặc ở $0^\circ$C.}

% ====== PHẦN IV ======
\section{PHẦN IV. TỰ LUẬN}
\resetcau

\cau Phân biệt sự khác nhau giữa biến đổi vật lí và biến đổi hóa học. Mỗi loại cho hai ví dụ minh họa.
\loigiai{\textbf{Biến đổi vật lí:} chất chỉ thay đổi hình dạng, kích thước hoặc trạng thái, KHÔNG tạo chất mới. Ví dụ: xé nhỏ tờ giấy; nước đá tan thành nước.\newline
\textbf{Biến đổi hóa học:} CÓ tạo thành chất mới (thường đổi màu, mùi, tính tan). Ví dụ: đốt giấy thành tro; đinh sắt bị gỉ.}

\cau a) Trình bày cách tách muối ăn ra khỏi dung dịch nước muối và giải thích vì sao làm được như vậy.\newline
b) Khi đun nóng hóa chất trong ống nghiệm, vì sao phải hướng miệng ống nghiệm về phía không có người?
\loigiai{\textbf{a)} Đun nóng dung dịch nước muối: nước bay hơi dần, muối không bay hơi nên kết tinh và ở lại, thu được muối. Làm được vì muối tan trong nước tạo dung dịch; khi nước bay hơi hết, muối tách ra ở dạng tinh thể.\newline
\textbf{b)} Vì khi đun, chất lỏng có thể sôi và bắn ra khỏi miệng ống; hướng miệng ống về phía không có người để tránh hóa chất nóng bắn vào người, bảo đảm an toàn.}

\bigskip
\begin{center}\rule{2cm}{0.6pt}\ \textbf{HẾT}\ \rule{2cm}{0.6pt}\end{center}

\end{document}
